\section{The exact logical strength visible at one native binder}
\label{sec:depth-one-boundary}

This section returns to the complete initialized-$I$ native calculus of
Section~\ref{sec:axioms}; its selector, predecessor, and PR compiler have
been established. Those resources have not been claimed for the weaker
self-seeded $F$ theory of the preceding section.

This observation locates a concrete boundary that a proposed simplification
of the induction guard must respect. It is a theorem about the established
arithmetic interpretation of the native syntax, not a non-bi-interpretability
theorem for every conceivable unrelated interpretation.

Write $C(0,a,b)=a$ and $C(Sn,a,b)=b$ for the derived numerical selector,
and $P(0)=0$, $P(Sn)=n$ for the derived predecessor. All formulas below may
have arbitrary fixed numerical parameters. A formula is called
$\Sigma_2^{\mathrm{qf}}$ here if it has the form
$\exists x\,\forall y\,\theta(x,y,\vec p)$ with quantifier-free $\theta$.
Several variables in a block may be packaged by the already derived
primitive-recursive pairing operations. This notation specifies an actual
prenex class; it does not appeal to classical prenex normalization of an
arbitrary intuitionistic formula.

\begin{proposition}[One-binder expressivity]
Under the whole-domain interpretation in iPRA, the class of native formulas
$D_{x,z}E(x,z,\vec p)$ with $E$ a literal equation is exactly the class
$\Sigma_2^{\mathrm{qf}}$, up to provable equivalence. Similarly, native
$Q_zE(z,\vec p)$ formulas are exactly the corresponding
$\Pi_1^{\mathrm{qf}}$ class, up to provable equivalence.
\end{proposition}

\begin{proof}
The forward direction follows from
\[
 D_{x,z}E\equiv
 \exists x\bigl(E(x,0)\to\forall z\,E(x,z)\bigr)
 \leftrightarrow
 \exists x\,\forall z\bigl(E(x,0)\to E(x,z)\bigr).
\]
The final matrix is quantifier-free. The displayed interchange of
implication and universal quantification is intuitionistically valid,
because $z$ is absent from $E(x,0)$. Dropping the outer existential gives
the analogous statement for $Q$.

Conversely, choose the iPRA characteristic term $r(x,y,\vec p)$ for a
given quantifier-free matrix, with
$r(x,y,\vec p)=0\leftrightarrow\theta(x,y,\vec p)$. Expand $C$, $P$, and $r$
into their fixed compiled initialized-iterator source terms, so the following
formula is literally an equation of the native language. Define the literal
equation
\[
 E(x,z,\vec p):\quad
 C\bigl(z,0,r(x,Pz,\vec p)\bigr)=0.
\]
Its zero instance is provable. Its successor instance is equivalent to
$r(x,y,\vec p)=0$, since $P(Sy)=y$. Numerical coverage therefore proves
\[
 \forall z\,E(x,z,\vec p)
 \leftrightarrow \forall y\,\theta(x,y,\vec p).
\]
Since $E(x,0,\vec p)$ is provable, its occurrence as the antecedent in
the native $D$ reading disappears by intuitionistic reasoning. Thus
$D_{x,z}E$ has precisely the required existential-universal meaning.
The same construction with no existential coordinate establishes the
$Q$ claim.
\end{proof}

\begin{corollary}[A one-binder progress extension already adds a full
prenex induction class]
Suppose the native equation-progress introduction rule is extended to
matrices $B(n,\vec p)$ that consist of one $D$ binder over a literal
equation. Over the existing native theory, this extension derives every
induction instance for $\Sigma_2^{\mathrm{qf}}$ formulas. Extending the
rule to one $Q$ binder over an equation derives the analogous
$\Pi_1^{\mathrm{qf}}$ induction instances.
\end{corollary}

\begin{proof}
Let $B(n,\vec p)$ be the native translation supplied by the proposition.
Suppose $\Gamma\vdash B(0,\vec p)$ and
$\Gamma,B(z,\vec p)\vdash B(Sz,\vec p)$, with the usual fresh $z$.
Apply the extended progress rule using an unused existential coordinate.
It yields $D_{u,n}B(n,\vec p)$. Paired elimination exposes
$Q_nB(n,\vec p)$, and closure elimination with the given base instance
proves $B(t,\vec p)$ for an arbitrary term $t$. The internal equivalence
from the proposition transports this derivation to the original prenex
formula. The argument for $Q$ matrices is identical.
\end{proof}

Consequently, retaining a test that merely says ``at most one binder''
does not preserve the stated literal-equation induction restriction.
Such a change supplies induction for this additional syntactic class;
equivalence with the original rule would require a separate elimination
or conservativity theorem. No claim here identifies every prenex
induction extension with full HA, and no strength separation for
arbitrary alternative bi-interpretations is inferred from this calculation.

\subsection{A uniform lower bound on what larger depths can express}

The one-binder equivalence extends to a useful inclusion at all depths.
The qualification is important: no intuitionistic prenex-normal-form
theorem is being silently used for nested native implications.
Let $\delta(E)=0$ for equations and
$\delta(Q_yA)=\delta(D_{x,y}A)=1+\delta(A)$.

For a fresh numerical parameter $z$, define $G_zA$ by replacing every
equation $s=t$ in the fixed finite formula $A$ by
\[
                 C(z,0,s)=C(z,0,t),
\]
and retaining its native binders. Capture is avoided by renaming the
bound variables first. This transformation does not change $\delta$.
It has the internal properties
\[
             \vdash G_0A,\qquad G_{Su}A\leftrightarrow A.
\tag{Gate}\label{eq:depth-gate}
\]
For equations these are computation identities. For a $Q$ or $D$ root,
the first assertion follows inductively because its entire matrix is
true, using the zero witness for $D$; the second follows from the same
interpreted binder applied to equivalent matrices. This reasoning uses
the already recovered full logical interface, rather than assuming a
new matrix-congruence rule during an earlier arithmetic bootstrap.

\begin{theorem}[Alternating-prefix encoding at depth $k$]
For every $k\geq0$, every explicitly prenex formula
\[
 \exists x_1\forall y_1\cdots\exists x_k\forall y_k\,
 \theta(\vec x,\vec y,\vec p),
\]
with quantifier-free $\theta$, has an iPRA-provably equivalent native
formula of depth at most $k$. The translation is uniform by external
recursion on the displayed finite prefix. In particular, allowing
equation-progress introduction on all matrices of depth at most $k$
derives induction for this entire prenex class.
\end{theorem}

\begin{proof}
For $k=0$, replace the quantifier-free formula by its primitive-recursive
characteristic equation. For the induction step suppose $F(x,y,\vec p)$
already translates the remaining prefix at depth at most $k-1$. Put
\[
             \mathcal B(F)
       :=D_{x,z}\,G_zF(x,Pz,\vec p).
\]
By \eqref{eq:depth-gate}, its zero matrix is true and its successor
matrix at $z=Sy$ is equivalent to $F(x,y,\vec p)$. Numerical coverage
then gives
\[
             \mathcal B(F)
       \leftrightarrow\exists x\,\forall y\,F(x,y,\vec p).
\]
Its depth is at most $k$, establishing the recursive construction.
For induction, repeat the derivation in the preceding corollary with
this translated matrix. The equivalences commute with the passive
induction parameter and its zero and successor substitutions.
\end{proof}

Blocks of several numerical variables can be packaged into one tuple
using the available pairing and projections. Dummy blocks can be added
because the numerical domain is inhabited. Consequently the inclusion
covers the usual explicitly prenex $\Sigma_{2k}$ class, and for $k\ge1$ the
$\Pi_{2k-1}$ class by an unused first existential block. It does not
claim that every intuitionistic formula of a conventional arithmetical
hierarchy has been prenexed without additional assumptions.

At depths greater than one, only this encoding inclusion is claimed
intuitionistically. The native reading contains implications whose
antecedents may themselves be quantified; moving such quantifiers
through implication is not a general intuitionistic equivalence.
Under classical logic the usual prenex manipulations give the familiar
$\Sigma_{2k}$ upper bound for a $D$ root of depth $k$, but that classical
observation must not be substituted for an iPRA proof. The new inclusion
and the exact depth-one calculation already suffice to audit the
literal-equation guard.

